%% kapa: exam=1 round=32 year=2021 period=2 form=A subject=accounting subject_ko=회계학 seq=32 number=72
%% source: Q-Net 감정평가사 기출문제 (cst00309) artlSeq=5208397 file=extracted/1008/5208397/A형/2021년도 제32회 감정평가사 1차 2교시 A형.pdf
%% conversion: kordoc 4.24.51 -> tools/gen1.py (auto); review=proofread-sonnet; flags=none
\begin{question}{72}
(주)감평은 단일 제품을 대량생산하고 있으며, 가중평균법을 적용하여 종합원가계산을 하고 있다. 직접재료는 공정초에 전량 투입되고, 전환원가는 공정 전체에서 균등하게 발생한다. 당기 원가계산 자료는 다음과 같다.
\begin{framedtext}
{\centering\begin{tabular}{@{\hspace{1em}}l@{\hspace{4em}}r@{}}
기초재공품 & 3,000개(완성도 80\%) \\
당기착수수량 & 14,000개 \\
당기완성품 & 13,000개 \\
기말재공품 & 2,500개(완성도 60\%)
\end{tabular}\par}
\end{framedtext}
품질검사는 완성도 70\%에서 이루어지며, 당기 중 검사를 통과한 합격품의 10\%를 정상공손으로 간주한다. 직접재료원가와 전환원가의 완성품환산량 단위당 원가는 각각 ￦30과 ￦20이다. 완성품에 배부되는 정상공손원가는?
\begin{choices}
  \item ￦35,000
  \item ￦44,000
  \item ￦55,400
  \item ￦57,200
  \item ￦66,000
\end{choices}
\end{question}
\answer{2}
